\chapter{Обработка результатов измерений}
\label{ch:measurement-processing}

\section{Где живёт метрология в серии}
\label{sec:metrology-layer}

Микроуровень~$\Mlayer$ в этой монографии не является объектом прямого прибора
(\cref{ch:macro}): доступен только макроскопический слой~$\Tlayer$ и корпуса
опубликованных констант (CODATA, PDG, астрофизические каталоги).
\textbf{Прямые измерения} относятся к величинам, которые протокол опыта
или международная шкала выдают как $x$ с заявленной неопределённостью.
\textbf{Косвенные} --- к любой величине $y=f(x_1,\ldots,x_n)$, собранной из прямых
(размерный анализ, QED-формула, SI-мост \cref{ch:si-sm}).

Сопоставление модели с опытом на~$\Tlayer$ \emph{обязано} следовать обычной
метрологической арифметике (рекомендации GUM, ISO/IEC Guide 98-3):
иначе «проценты совпадения'' не интерпретируются.
Глава фиксирует единый язык анализа для \cref{ch:computational-laboratory}
и развёрнутых записей в «Критических экспериментах'' каждого тома.

\begin{remark}[Теорема на $M$ vs число на $T$]
Тождество на~$\Mlayer$ (например, замыкание $g$) не несёт экспериментального $u(x)$,
пока не прочитано на~$\Tlayer$ через осреднение и прибор.
Сравнение «модель $\approx$ PDG'' всегда относится к \emph{прочитанному} $y$,
не к сырому состоянию решётки.
\end{remark}

\section{Прямые измерения}
\label{sec:direct-measurement}

\begin{definition}[Результат прямого измерения]
Величина $X$ измерена $n$ раз в воспроизводимых условиях;
отчётный результат $\hat x$ и \emph{стандартная неопределённость} $u(x)$
описывают наилучшую оценку и разброс, согласованный с протоколом
(тип~A --- по статистике повторов; тип~B --- по модели прибора, калибровке, хвостам распределения).
\end{definition}

Расширенная неопределённость для уровня охвата примерно 95\,\%:
\begin{equation}
\label{eq:expanded-uncertainty}
  U(x) = k\,u(x),\qquad k\approx 2
\end{equation}
при нормальном эффективном распределении и выборе $k$ под заявленный охват.
В корпусах констант $u(x)$ уже свёрнута в табличное значение; в тексте серии
$\hat x$ --- табличный центр, $u(x)$ --- его стандартная компонента, если явно не оговорено иное.

\section{Косвенные измерения и перенос неопределённости}
\label{sec:indirect-gum}

Пусть
\begin{equation}
  y = f(x_1,\ldots,x_n).
\end{equation}
При малых отклонениях от наилучших оценок $\hat x_i$ (линейный перенос, GUM):
\begin{equation}
\label{eq:gum-linear}
  u^2(y) = \sum_{i=1}^{n}\left(\frac{\partial f}{\partial x_i}\right)^{\!2} u^2(x_i)
  + 2\sum_{i<j}\frac{\partial f}{\partial x_i}\frac{\partial f}{\partial x_j}\,r_{ij}\,u(x_i)\,u(x_j),
\end{equation}
где $r_{ij}$ --- коэффициент корреляции входов (если входы независимы, второй ряд обнуляется).

Для произведения степеней $y = c\,x_1^{a_1}\cdots x_n^{a_n}$ удобна относительная форма:
\begin{equation}
\label{eq:gum-relative}
  \left(\frac{u(y)}{|y|}\right)^{\!2}
  \approx \sum_i a_i^2\left(\frac{u(x_i)}{x_i}\right)^{\!2}
  + 2\sum_{i<j} a_i a_j r_{ij}\frac{u(x_i)}{x_i}\frac{u(x_j)}{x_j}.
\end{equation}

\begin{remark}[Пример: ведущий порядок QED для пара-позитрония]
\label{ex:para-ps-tau}
Время жизни $^1S_0$:
\begin{equation}
\label{eq:para-ps-tau}
  \tau = \frac{2\hbar}{m_e c^2\,\alpha_{\mathrm{fs}}^{\,5}}.
\end{equation}
Из \eqref{eq:gum-relative} при независимых $m_e$ и $\alpha$:
\begin{equation}
\label{eq:tau-relative-unc}
  \frac{u(\tau)}{\tau}
  \approx \frac{u(m_e)}{m_e} + 5\,\frac{u(\alpha_{\mathrm{fs}})}{\alpha_{\mathrm{fs}}}.
\end{equation}
Систематика \emph{модели} (пропущенные радиационные поправки к $\Gamma$)
добавляется отдельно как тип~B и не подменяет $u(m_e)$ из таблицы.
\end{remark}

\section{Выборка, стандартная ошибка и Стьюдент}
\label{sec:student}

Повторы опыта (или независимые прогоны симуляции-прокси) дают выборку $x_1,\ldots,x_n$.

\begin{definition}[Выборочные оценки]
\label{def:sample-stats}
Наилучшая оценка среднего $\bar x = \frac{1}{n}\sum_i x_i$.
Несмещённая выборочная дисперсия
$s^2 = \frac{1}{n-1}\sum_i (x_i-\bar x)^2$,
стандартная ошибка среднего $\mathrm{SE}(\bar x) = s/\sqrt{n}$.
\end{definition}

При нормальности $x_i$ интервал для математического ожидания при охвате $1-\alpha$:
\begin{equation}
\label{eq:student-ci}
  \bar x \pm t_{n-1,\,1-\alpha/2}\,\frac{s}{\sqrt{n}},
\end{equation}
где $t_{n-1,\,1-\alpha/2}$ --- квантиль распределения Стьюдента с $n-1$ степенями свободы.
При $n\gtrsim 30$ часто используют $k\approx 2$ вместо $t$, но для малых ансамблей
(десятки и меньше) замена $t\to 2$ завышает уверенность в совпадении.

\begin{remark}[Симуляция и тип~A]
Статистика по ансамблю начальных условий измеряет разброс \emph{прокси-прибора}
на решётке, а не полную неопределённость PDG.
Интервал \eqref{eq:student-ci} годится для «насколько стабильна ось / пик / задержка'',
но не отменяет тип~B для «насколько прокси соответствует bound-state $\tau$''.
\end{remark}

\section{Сравнение с эталоном}
\label{sec:compare-reference}

Пусть модель (или косвенная цепочка на~$\Tlayer$) даёт $y_{\mathrm{mod}}$,
эталон --- $y_{\mathrm{ref}}\pm u_{\mathrm{ref}}$ (PDG, CODATA).
При независимых источниках неопределённости:
\begin{equation}
\label{eq:combined-unc}
  u_{\mathrm{comb}} = \sqrt{u_{\mathrm{mod}}^2 + u_{\mathrm{ref}}^2},
\end{equation}
а согласие проверяют неизвестностью
\begin{equation}
\label{eq:en-score}
  E_n = \frac{|y_{\mathrm{mod}}-y_{\mathrm{ref}}|}{u_{\mathrm{comb}}}.
\end{equation}
Правило $E_n \lesssim 1$ (или $2$ при охвате 95\,\% и корректном $k$) --- \emph{метрологическое},
не замена физического вывода: при заведомо неполной модели
(ведущий порядок без петлевых поправок) допустимый остаток задаётся
явной систематикой типа~B, а не сужением $u_{\mathrm{ref}}$.

\begin{remark}[Круговость входа]
Если $y_{\mathrm{mod}}$ собран из тех же калиброванных $x_i$, что уже вошли
в определение $y_{\mathrm{ref}}$ (типично: CODATA $m_e$ в формуле и PDG-окно,
натянутое на ту же шкалу), сравнение становится \emph{косметическим}.
Честная проверка требует независимого входа (другая лестница масс, upstream-оценка
из замкнутой теории без подстановки того же табличного $m_e$) или явного учёта корреляции.
\end{remark}

\section{Несколько наблюдаемых}
\label{sec:multiple-observables}

Одна и та же безразмерная константа ($\alpha_{\mathrm{fs}}$, $N_c$)
входит в несколько косвенных $y_j$: совпадения не независимы.
Для набора $(y_1,\ldots,y_m)$ с ковариационной матрицей $\Sigma$
естественна метрика
\begin{equation}
  \chi^2 = ({\mathbf y}_{\mathrm{mod}}-{\mathbf y}_{\mathrm{ref}})^{\mathsf T}
  \Sigma^{-1} ({\mathbf y}_{\mathrm{mod}}-{\mathbf y}_{\mathrm{ref}}),
\end{equation}
с $\chi^2/(m-p)$ близким к единице при $m-p$ степенях свободы и $p$ действительно свободных параметрах.
В монографии достаточно \emph{не считать} пять близких SI-якорей пятью независимыми «победами'';
в критических экспериментах указывать зависимые наблюдаемые явно.

\section{Связь с критическими экспериментами}
\label{sec:metrology-critical}

Блок «Критические эксперименты'' в каждом томе формулирует предсказание и критерий провала.
Обязательные поля согласования с настоящей главой:
\begin{enumerate}[label=\textbf{M\arabic*:},leftmargin=*]
\item класс величины: прямая / косвенная / sim-прокси;
\item откуда $u$ (таблица, \eqref{eq:gum-linear}, ансамбль + \eqref{eq:student-ci});
\item отделена ли систематика модели (тип~B) от статистики повторов (тип~A);
\item нет ли кругового входа (\cref{sec:compare-reference});
\item для набора якорей --- зависимость через общие $x_i$ (\cref{sec:multiple-observables}).
\end{enumerate}

Внутренний реестр согласованности кода (verify) контролирует алгебру и воспроизводимость;
он не подменяет протокол $E_n$ относительно внешнего эталона, если только verify
не реализует явно \eqref{eq:gum-linear}--\eqref{eq:en-score} с заявленными $u$.
